\input{preamble}
\title{More than One Way to Skin a Cat:\\
Preserving Verified Modes in RLVR}
\author{Anonymous Authors}
\begin{document}
\maketitle

\begin{abstract}
Mathematical agents need candidates for search, verification, and recovery.
Yet training can improve correctness while narrowing this repertoire. We
introduce \mb{}, which measures distinct verified outputs across five executable
domains. ReplayMaxRL combines sampling-aware reinforcement learning with memory
of verified outputs, improving correctness while preserving alternatives.
Replay helps under Dr.GRPO and MaxRL across three and two model scales
respectively, and on a harder graph-coloring task. These results separate
discovering successful behaviors from preserving them.
\end{abstract}

\section{Preserving alternatives for mathematical agents}
\label{sec:introduction}\label{sec:collapse}
Mathematical agents propose candidates, check them with tools, and revise after
feedback \citep{brown2024monkeys,snell2024testtime}. Alternative proof or
computation routes could aid search and self-correction when attempts fail or
constraints change. A verifier can only select from what the policy still proposes.

Reinforcement learning with verifiable rewards (RLVR) can improve correctness
while these alternatives disappear \citep{guo2025deepseekr1}
(Figure~\ref{fig:story}).

\begin{figure}[!htbp]
\centering
\includegraphics[width=\linewidth]{figures/modecollapse_story.pdf}
\caption{\textbf{Correctness can rise while alternatives disappear.}
One graph admits six valid colorings. Bars pool four eight-sample draws from
matched Qwen2.5-3B runs. At pass 6.5, Dr.GRPO produces 22/32 correct from one
mode; ReplayMaxRL produces 32/32 across four.}
\label{fig:story}
\end{figure}

The loss also appears across our datasets. Averaged over five domains, eight
passes of Dr.GRPO reduce the initial number of verified alternatives beyond
the first correct output in eight samples by $99.4\%$ at Qwen2.5-0.5B,
$71.6\%$ at Falcon3-1B, and $41.5\%$ at Qwen2.5-3B
(Appendix Figure~\ref{fig:baseline-precheck}).

Our idealized analysis links retention at matched accuracy to updates shared
across correct outputs (Appendix~\ref{app:theory-size},
Theorem~\ref{thm:shared-correctness}); size alone provides no guarantee.
We ask how training can retain alternatives after they become rare.

\section{ModeBench makes output identity executable}
\label{sec:modebench}
To distinguish a different solution from a different string, \mb{} assigns each
accepted output a canonical key; each distinct key defines a verified
\emph{mode}. The validator returns both correctness and identity
(Figure~\ref{fig:modebench-examples}): computation routes in
Countdown and MathIR, tested program behaviors in Python, and valid
configurations in Graph and PantryPlan. Cosmetic aliases merge; distinct routes can share an answer
(key rules: Appendix~\ref{app:domain-overview}).

\begin{figure}[!htbp]
\centering
\includegraphics[trim={0 9bp 0 0},clip,width=\linewidth]{figures/modebench_examples.pdf}
\caption{\textbf{One prompt, different verified outputs.}
Keys retain a color vector, normalized computation tree, tested return vector,
algebraic state trajectory, or ingredient support. Cosmetic aliases merge;
these keys identify task-defined outputs rather than free-form reasoning.}
\label{fig:modebench-examples}
\end{figure}

\section{ReplayMaxRL: learn from new samples and remember successes}
\label{sec:method}\label{sec:maxrl-method}\label{sec:replay}
These keys let us store distinct verified solutions and revisit them when
they are absent from fresh rollouts. ReplayMaxRL combines these replay
updates with MaxRL learning from new samples
(Figure~\ref{fig:verified-support-story}).

\begin{figure}[!htbp]
\centering
% Crop only blank PDF margins; keep the figure itself unchanged.
\includegraphics[trim={4bp 23bp 4bp 10bp},clip,width=\linewidth]{figures/verified_support_story.pdf}
\caption{\textbf{Discovery and preservation use the same verified samples.}
MaxRL learns from fresh rewards; replay retains exemplars by key and revisits
rare successes. Memory cannot protect a mode before it is discovered.}
\label{fig:verified-support-story}
\end{figure}

MaxRL trains policy $\pi_\theta$ for success within $k$ samples across sampling
budgets \citep{tajwar2026maxrl}. Each training prompt $x$ keeps one
policy-generated exemplar $e_c$ per admitted key in bank $\mathcal B_x$.
With $|e_c|$ its token length, recurrent replay minimizes
\begin{equation}
L_{\mathrm{replay}}(x)=-\frac{1}{|\mathcal B_x|}
\sum_{c\in\mathcal B_x}\frac{\log\pi_\theta(e_c\mid x)}{|e_c|}.
\label{eq:lm-verified-replay}
\end{equation}
Equal averaging of per-token scores gives each retained key equal weight. Replacing MaxRL with Dr.GRPO \citep{liu2025understanding},
which debiases GRPO \citep{shao2024deepseekmath}, yields our ReplayDr.GRPO.
Appendix~\ref{app:algorithm} gives
the full update. Appendix~\ref{app:theory} proves categorical collapse and
replay retention (Theorems~\ref{thm:grpo-collapse} and~\ref{thm:replay-retention}).

\noindent\textbf{Related work.} Prior work studies diversity collapse
\citep{cui2025entropy,dang2025diversity}, sampling-aware objectives, and success
replay \citep{zhang2025rlep}. Our focus is execution-defined identity for
measuring and retaining verified alternatives. Appendix~\ref{sec:related}
gives the full discussion.

\section{Does memory preserve useful output support?}
\label{sec:experiments}\label{sec:results}
We test whether remembering training successes benefits new problems.
\texttt{pass@8} is the chance that eight samples contain any correct response;
\texttt{distinct@8} is their average number of distinct verified modes. Reporting
both separates success from the alternatives available for selection, although
mode counts remain coupled to correctness (definitions: Appendix~\ref{app:metrics}).

Five seeds are registered per comparison for eight passes; each of 128 held-out prompts per
domain receives four eight-sample draws. Controls execute the replay code with its gradient zeroed. Figure~\ref{fig:cross-scale-terminal-effects} compares
Dr.GRPO with replay and alternatives targeting reward normalization, on-policy
diversity (UCPO; \citealp{lochab2026ucpo}), success replay (RLEP-Dr),
and semantic entropy. Protocols and uncertainty are in
Appendices~\ref{app:experimental-design} and~\ref{app:extended-results}.

\begin{figure}[!htbp]
\centering
\includegraphics[width=\linewidth]{figures/experiment1_retention_comparator_matrix.pdf}
\caption{\textbf{Replay improves correctness and verified support across scales.}
(A) ReplayDr.GRPO minus Dr.GRPO. (B) Qwen2.5-0.5B alternatives and the untrained
reference minus Dr.GRPO. Daggers mark four admissible pairs; others use five. Average uses shared seeds.
Left: probability differences. Right: key counts. Bold marks column maxima;
uncertainty: Appendix~\ref{app:supporting}.}
\label{fig:cross-scale-terminal-effects}
\end{figure}

\noindent\textbf{Replaying stored solutions improves candidates for new problems.}
\label{sec:results-retention}
Replay revisits a bank of verified solutions discovered during training.
ReplayDr.GRPO improves both means in all 14 complete blocks and the remaining
four-seed block; \texttt{distinct@8} rises in all 74 admissible pairs
(source exclusion: Appendix~\ref{app:source-integrity}). On Qwen2.5-0.5B Countdown,
\texttt{pass@8} rises from $.48$ to $.67$ and \texttt{distinct@8} from $.49$
to $1.64$. The gain includes more computation routes alongside greater success.
These held-out prompts never enter the bank: replay improves candidates for new problems.
Across five Qwen2.5-0.5B domains, uniform versus frequency-weighted replay adds
$.318$ extra modes on average (95\% bootstrap interval $[.272,.358]$;
Appendix~\ref{app:frequency-replay-progress}).

\clearpage
\begin{figure}[!htbp]
\centering
\includegraphics[width=\linewidth]{figures/e118_all_scale_factorial_progress.pdf}
\caption{\textbf{Memory adds value beyond sampling-aware MaxRL.}
Each track connects the untrained model, fresh-objective training, and its replay
counterpart. Faint paths are seeds; bold paths average domains within seed.
Falcon Dr.GRPO uses four common admissible seeds; other tracks use five. The
Qwen2.5-3B row shows its completed Dr.GRPO/replay track.}
\label{fig:maxrl-factorial}
\end{figure}

\noindent\textbf{Stronger fresh learning still benefits from memory.}
\label{sec:results-maxrl}
We cross Dr.GRPO/MaxRL with replay absent/present
(Figure~\ref{fig:maxrl-factorial}). Replay improves both cross-domain averages
under both objectives at the two completed factorial scales. Qwen2.5-3B also
benefits on its completed Dr.GRPO track; Falcon Dr.GRPO is descriptive at four
seeds. For ReplayMaxRL, both
endpoints improve in all five Qwen2.5-0.5B domains; Falcon correctness improves in five
and breadth in four. Cross-domain averages are post-hoc summaries; paired intervals
are in Appendix~\ref{app:results-maxrl}. Optimizing success across sampling
budgets still benefits from revisiting earlier successes.

\begin{figure}[!htbp]
\centering
\includegraphics[width=\linewidth]{figures/modebench_level_admission.pdf}
\caption{\textbf{Replay across difficulty levels with Qwen2.5-0.5B-Instruct.}
Open/filled markers: Level 1/2. (A) Frozen-model success.
(B) Interim means weight the five domains equally after averaging seeds.
Checkpoints match across levels and methods; Pantry has one seed at step 0
(Appendix~\ref{app:level2-interim}).}
\label{fig:level2-admission}
\end{figure}

\noindent\textbf{Replaying stored solutions also helps on harder problems.}
\label{sec:results-levels}
Level 2 increases difficulty while preserving the validator, key rule, and
valid-mode-count histogram (Appendix~\ref{app:data}). Replay raises both interim
means at both levels in Figure~\ref{fig:level2-admission}B. The completed Graph
block provides an eight-pass example: Dr.GRPO/MaxRL reach $.20/.43$
\texttt{pass@8}, versus $.76/.74$ with replay.

\section{Conclusion and Limits}
\label{sec:conclusion}\label{sec:discussion}\label{sec:limitations}
ModeBench makes a missing part of evaluation measurable: the range of distinct,
verified outcomes a model can produce for one problem. This matters because
different approaches can have different strengths. One may succeed where
another fails, or remain useful when constraints change. Understanding diversity
therefore gives a fuller picture of a model's capabilities than accuracy alone.

Our results show that replay can preserve alternatives while improving
held-out correctness. These controlled tasks provide a starting point for
studying diversity in broader mathematical reasoning. Replay retains only discovered,
stored modes; individual-mode survival remains unmeasured
(Appendix~\ref{app:ablations}). We have not shown that broader support causes
accuracy gains or helps solve later problems. The next step is to test which
alternatives help on richer mathematical tasks, and when.

\label{main:lastpage}
\clearpage
\bibliographystyle{plainnat}
\bibliography{example_paper}
\clearpage
\appendix
\input{appendix}
\end{document}
